\documentclass[addpoints]{exam}
% \documentclass[addpoints,12pt]{exam}

\usepackage[slovene]{babel}
\usepackage{amsfonts}
\usepackage{amssymb}
\usepackage{amsmath}
\usepackage{float}
\usepackage{graphicx}
\usepackage{enumitem}
\usepackage{color}
\usepackage{eurosym}

\usepackage{pgfpages}
% \usepackage{tikz}
\usepackage{wrapfig}
\usepackage{pgfkeys}
\usepackage{pgfplots}
\usepackage{xcolor}
\usepackage{tkz-euclide}
\usepackage{xfp}




%Komentar naslednje vrstice glede na verzijo testa -- rešitve ali prostor za odgovore:
% \printanswers

% \linespread{2}


\pointsinrightmargin
\marginbonuspointname{}


\renewcommand{\solutiontitle}{\noindent\textbf{Rešitve in točkovanje:}\par\noindent}

\colorgrids
\definecolor{GridColor}{gray}{0.7}
\setlength{\gridsize}{7mm}

\extrawidth{5mm}
\setlength{\rightpointsmargin}{1.5cm}

\begin{document}

\runningfooter{}{}{}

\begin{center}
    \fbox{\fbox{\parbox{15cm}{\centering
    Popravljanje in pridobivanje ocen -- peto pisno ocenjevanje znanja \\ 2. letnik -- test C \\ 15. junij 2026 \\ (Kvadratna funkcija; Kompleksna števila)}}}
\end{center}

\vspace{0.1in}
\makebox[\textwidth]{Ime in priimek, oddelek:\enspace\hrulefill}


\begin{center}
    \hqword{Naloga}
    \hpword{Možne točke}
    \chbpword{Dodatne točke}
    \hsword{Dosežene točke}
    \htword{Skupaj}  
    % \combinedpointtable[h][questions]
    \gradetable[h][questions]

\end{center}

\vspace{0.1in}


\begin{questions}

    % 1. vprašanje
    \question[6]
        Dani sta družina parabol $y=(2a+1)x^2-12x+9$ in parabola $y=-x^2+2x-3$.
        Kateremu pogoju mora ustrezati parameter $a$, da se paraboli dotikata?
        Določite enačbo parabole.

            \begin{solution}[8cm]

            \end{solution}


    % 2. vprašanje
    \question
        Miha vrže z okna svoje sobe v zrak kovanec. Funkcija $h$, s predpisom $h(t)=-5t^2+10t+15$, 
        opisuje višino kovanca v odvisnosti od časa. Čas je merjen v sekundah, višina pa v metrih.
        
        \begin{parts}
            \part[1] S katere višine je Miha vrgel kovanec?
            \part[3] Po kolikšnem času pade kovanec na tla?
            \part[2] Kolikšna je največja višina kovanca med letom?
        \end{parts}

            \begin{solution}[3cm]

            \end{solution}


            \newpage
    % 3. vprašanje
    \question
        Dano je kompleksno število $$z= \left(3-i\right)^4-i^{39}-\left(2+i\sqrt{3}\right)\left(2-i\sqrt{3}\right)+\dfrac{5+3i}{3-5i}. $$

        \begin{parts}
            
            \part[7]
            Poenostavite zapis kompleksnega števila $z$.

                \begin{solution}[14cm]

                \end{solution}

                
            \part[1]
            Določite realni del kompleksnega števila $z$.

                \begin{solution}[2cm]

                \end{solution}


            
            \part[1]
            Določite imaginarno komponento $Im(z)$ kompleksnega števila $z$.

                \begin{solution}[2cm]

                \end{solution}


            \part[1]
            Določite imaginarno del kompleksnega števila $z$.

                % \begin{solution}[2cm]

                % \end{solution}
                
        \end{parts}


    \newpage

    % 4. vprašanje
    \question[6]
        V kompleksni ravnini upodobite množico točk:                     
            $$\left\{z\in\mathbb{C}; \left|Re(z)-1\right|>2 \land \left|z\right|\leq 4 \right\}. $$ 
        
            \begin{solution}[12cm]

            \end{solution}


    % 5. vprašanje
    \question[6]
        V množici kompleksnih števil rešite enačbo: $$x^4-3x^2=4.$$

            \begin{solution}[8cm]

            \end{solution}









\end{questions}



\end{document}